\documentclass{letter}

% Top margin, right margin, left margin, bottom margin, footnote skip
\usepackage[tmargin=3.5cm,rmargin=3cm,lmargin=3cm,bmargin=3cm]{geometry} 
\usepackage[utf8]{inputenc}
\usepackage{biblatex}
\addbibresource{./references.bib}
% linktocpage shall be added to snippets.
\usepackage{hyperref,theoremref}

\usepackage{setspace}
\usepackage{enumitem}

\linespread{1.2}

\address{2 Winchester Rd \\ Oxford}
\name{Me}
\signature{My Signature line}

\title{Application to Stipendiary Lectureship in Computer Science}
\author{Yuhao Han} 
\date{\today}

\begin{document}
% \tableofcontents

\begin{letter}{Jesus College \\ University of Oxford}
\opening{Dear Members of the Selection Committee,}

I am writing to apply for the Stipendiary Lectureship in Computer Science at Jesus College.
I am currently completing my Master's studies in Mathematics and Foundations of Computer Science at Oxford and will begin my DPhil studies in quantum computing in Michaelmas Term 2026.

This position appears appealing to me, as I am passionate about teaching and it would allow me to contribute to the social and intellectual life of the College. 
I am excited to share my knowledge and grow together with my students through small-group tutorial settings.

I have a strong educational background that spans computer science, programming, and mathematics.
During my undergraduate and Master's studies, I have consistently achieved excellent marks above 80.
I have research experience on the Linux kernel, compilers, AI, and algebraic geometry. 
I have also developed substantial software and am familiar with several programming languages, including Python, C++, and Rust.
My mathematical training includes linear algebra, homological algebra, category theory, topology, among other subjects.
These experiences qualify me to teach various undergraduate courses in computer science including but not restricted to algorithms, imperative and functional programming, and linear algebra. I am also willing to undertake any other duties associated with this role, including providing pastoral care, marking, and helping with admissions and College events.

Since my undergraduate studies, I have engaged in various teaching and mentoring activities. 
At the University of Edinburgh, I volunteered for two years as a tutor teaching linear algebra to first-year students and as a mentor organising activities and providing pastoral care for junior students.
I have worked as an English tutor in China, teaching both adults and high-achieving school pupils.
Through these experiences, I have gained important perspectives on the academic and psychological needs of students from different backgrounds and with different levels of confidence.
They have also helped me develop my communication skills and my ability to explain complex ideas in accessible ways.

Thank you very much for considering my application.
Professor Aleks Kissinger at aleks.kissinger@cs.ox.ac.uk and Professor Sergii Strelchuk at sergii.strelchuk@cs.ox.ac.uk have agreed to be my referees.
I look forward to discussing my application further with you in the future.

Yours faithfully,\\
Yuhao Han

\newpage

\textbf{Teaching Statement}

Throughout my teaching career, I have adhered to the following two teaching principles:
\begin{enumerate}
   \item Teaching is more effective with ample examples, practice, and exercises.
   \item Each student has unique talents and perspectives, and teaching should be adjusted to their needs.
\end{enumerate}

During teaching, I will present to students ample examples and exercises. I will encourage them to complete the exercises and guide them towards understanding if they are stuck.
I will not give way answers or let them simply memorise facts.
I will also actively ask students for their feedback and adjust my teaching methods accordingly.

For the subject of computer science in particular, I believe practice is as important as, if not more important than, theory.
Therefore, I will encourage students to implement the theories they have learned in actual software on their computers.
I will also introduce students to practical programming tools, such as Git, Bash, and SSH; encourage them to try different ways of programming, such as using the command line and Linux; and teach them industry standards in software engineering.

In recent years, generative AI has posed both opportunities and challenges in education. 
I will follow the University's and the College's statements on AI.
I will encourage students to use AI to aid their studies when appropriate and also remind them to prioritise their own skills in coding, writing, and understanding complex concepts.
I will, in particular, warn them not to rely on AI for studying, as AI hallucinations, or the fabriacation of false facts, could be extremely detrimental.

Lastly, I believe teaching is a mutually beneficial activity, as expressed by the Confucian saying \textit{studying and teaching complement each other}, as it exposes me to diverse perspectives and helps me gain a better understanding of the subject.

\textbf{Research Interests}

My interests in mathematics and computer science stem from the following firm belief: mathematics, while itself beautiful and attractive, can, when applied correctly, create a great positive impact on our society through computer science.
The prime evidence of this is the internet and artificial intelligence, both of which have their foundations in advanced mathematics.
For my current master's dissertation, I am exploring a new tensor-contraction algorithm combining graph theory and ZX calculus, which is a categorical representation of linear maps as visually intuitive graphs.
I am actively exploring the mathematical theory of contraction algorithms, as well as implementing the algorithm in Rust.

My DPhil project will be a theoretical investigation of quantum circuit complexity using geometric methods. It studies how quantum operations can be modelled as paths in spaces of unitary transformations, with the aim of understanding the resources needed to implement these operations using standard quantum gates. 
The research will analyse reduced and symmetry-constrained models, including few-qubit systems and quotients by permutation symmetries or the Clifford group, in order to make the mathematical problem more tractable. 
The work is purely mathematical and computational, using differential geometry, Lie groups, Clifford algebras, quantum information theory, and symbolic and numerical computation. 

These research projects are directly related to undergraduate teaching, as they involve advanced ideas in linear algebra, algorithms, and proof systems. 
I would also introduce my research to students to help them understand the frontiers of research.

\end{letter}
\end{document}
